\silentchapter{Propositions}{propositions}

\begin{widebody}
\begin{center}\footnotesize
\begin{tabular}{@{}lllll@{}}
\textbf{\textit{d}} definition & \textbf{\textit{a}} postulate & \textbf{\textit{t}} theorem & \textbf{\textit{o}} observation & \textbf{\textit{i}} illustration \\
\end{tabular}
\end{center}
\end{widebody}

\vspace{6mm}

\mainprop{1}{}{The world of forms is all that is the case.}

\prop{1.1}{o}{The world of forms is the totality of realised configurations, not of things.}

\prop{1.11}{t}{Each realised configuration is specific; no configuration is partially realised.}

\prop{1.12}{t}{That a configuration takes an exact form, when infinitely many others were logically possible, demands an explanation beyond the configuration itself.}

\prop{1.13}{t}{The explanation lies in the relation between the manifested configuration and the total set of configurations that were possible but not realised.}

\prop{1.131}{t}{A form without alternatives is not a form; it is a necessity.}

\prop{1.2}{d}{An \textbf{object} is the simplest constituent of a configuration. The object is indivisible and immutable.}

\prop{1.201}{t}{The partition into objects is a choice of level of analysis. Below the chosen level physics continues; but not form-giving.}

\prop{1.203}{t}{The objects constitute the substance of the world of forms. Substance is what exists independently of what is the case.}

\prop{1.2031}{t}{If the world of forms had no substance, the question of whether a configuration has meaning would depend on whether another configuration is realised.}

\prop{1.204}{t}{The object is the immutable; the configuration is the unstable.}

\prop{1.21}{d}{A \textbf{configuration} is a determinate combination of objects. Form is the structure of the configuration.}

\prop{1.211}{i}{In the configuration the objects hang together like the links of a chain.}

\prop{1.212}{t}{The way the objects hang together in the configuration is the configuration's structure.}

\prop{1.213}{t}{Form is the possibility of structure.}

\prop{1.214}{t}{Two configurations with the same objects but different structure are different forms.}

\prop{1.2141}{t}{Two configurations with different objects but the same structure share form.}

\prop{1.22}{t}{Since it lies in the nature of the objects to be able to enter into configurations, all possible forms are already given with the objects.}

\prop{1.221}{t}{The possibility of entering into a configuration is the object's form.}

\prop{1.222}{t}{The morphospace is not constructed. It is discovered.}

\prop{1.2221}{t}{The designer does not invent possibilities; she uncovers them.}

\prop{1.223}{t}{To know the object is to know all its possible configurations. A new possibility cannot arise later.}

\prop{1.3}{d}{The \textbf{morphospace} of a class is the set of all possible configurations for objects in that class.\textsuperscript{2,\,10,\,20}}

\prop{1.31}{d}{Each realised object is one exact point in this n-dimensional space. Each axis is a measurable property.}

\prop{1.311}{t}{The choice of axes is a projection. Form exists independently of the measurement, as a mountain exists independently of the map.}

\prop{1.32}{d}{The empirical morphospace is always a projection. No finite set of parameters captures the latent complexity in full.\textsuperscript{2}}

\prop{1.321}{t}{That the morphospace is a projection does not mean that it is arbitrary. Some projections capture more variation than others. The best projection is the one that makes the most selection pressures visible.}

\prop{1.4}{d}{The morphospace partitions topologically into three regions: the \textit{settled}, the \textit{open}, and the \textit{forbidden}.}

\prop{1.41}{o}{The settled regions constitute the historical archive. They function as anchor points for all future navigation and constitute the inductive premise for further form.}

\prop{1.411}{t}{The weight of a settled position grows with occupation time. Independent convergence from isolated traditions confirms that the attractor is real, not conventional.}

\prop{1.412}{t}{The totality of settled positions is the history of form. The history of form determines what is known to be possible.}

\prop{1.42}{o}{The open regions are the \textit{adjacent possible}. Accessible but unactualised. They are finite; stretched between the settled and the forbidden.}

\prop{1.421}{t}{An open region bordering on a settled region is accessible. An open region bordering only on forbidden regions is inaccessible without a material change.}

\prop{1.43}{d}{The forbidden regions are excluded by material, structural, or energetic conditions. The prohibition follows the conditions, not the region. A material change can open a forbidden region.}

\prop{1.431}{t}{The boundary between open and forbidden is material, not logical. Beech bends along the grain; granite does not. The material decides what is forbidden.}

\prop{1.432}{t}{That the forbidden is materially bounded means that the world of forms can expand. Each new material knowledge shifts the boundary between the open and the forbidden.}

\prop{1.433}{t}{The topological regions are not identical to the epistemic: a forbidden area may be conceivable (materially bounded); an open region may be inconceivable (cognitively inaccessible).}

\prop{1.44}{d}{The three regions constitute an epistemic partition.\textsuperscript{16,\,17} The \textbf{known} (K), the \textbf{conceivable} (C\msym{∖}K), and the inconceivable. The boundaries are epistemic, not ontological.}

\prop{1.441}{d}{C is the \textit{conception space}. All forms that are conceivable but not settled. K is the \textit{knowledge space}. All forms known to work or to fail. Form-making is the movement between them.}

\prop{1.442}{t}{The totality of realised configurations determines not only what is known to be possible, but also what is known to be impossible. K contains both successes and failures.}

\prop{1.443}{t}{The boundary between C and the inconceivable is invisible from within. One cannot know what one cannot think.}

\prop{1.444}{t}{The operations between C and K span four directions:}

\begin{center}\small
\begin{tabular*}{\textwidth}{@{\extracolsep{\fill}}p{32mm}lp{22mm}@{}}
\toprule
\textbf{C-K operator} & \textbf{Direction} & \textbf{Light-cone operation} \\
\midrule
\textbf{C \msym{→} C}\newline{\footnotesize(concept expansion)} & Within C & Extension \\[4pt]
\textbf{C \msym{→} K}\newline{\footnotesize(conjecture \msym{→} knowledge)} & From C to K & Collapse \\[4pt]
\textbf{K \msym{→} C}\newline{\footnotesize(knowledge \msym{→} conjecture)} & From K to C & Motion \\[4pt]
\textbf{K \msym{→} K}\newline{\footnotesize(knowledge refinement)} & Within K & Resolution increase \\
\bottomrule
\end{tabular*}
\end{center}

\flowchapter{2 Not all positions are equally probable}{not all positions are equally probable}

\mainprop{2}{}{Not all positions are equally probable.}

\prop{2.01}{t}{That not all positions are equally probable means that something acts on the distribution. What acts is selection pressure.}

\prop{2.011}{t}{The pressure is not a force. It is a constraint on where the form \textit{cannot} be.}

\prop{2.1}{d}{A \textbf{selection pressure} modifies the probability distribution in the morphospace.\textsuperscript{3} The pressure is a gradient, not a rule.}

\prop{2.101}{t}{Form is what remains once the pressures have excluded everything else.\textsuperscript{15,\,77,\,78}}

\prop{2.11}{t}{A pressure that excludes nothing is no pressure. A pressure that excludes all but a single region is determinism. All observed pressures lie between these.}

\prop{2.111}{t}{Between the null pressure and determinism lies the entire space for design.}

\prop{2.12}{a}{Several selection pressures act simultaneously.}

\prop{2.121}{t}{With only one pressure, all forms would be identical.}

\prop{2.122}{t}{A latent pressure still acts; demonstrating it requires variation in the conditions.}

\prop{2.13}{a}{At least two selection pressures point in contradictory directions. Each realised form is therefore a compromise.}

\prop{2.131}{t}{The term ``suboptimal'' presupposes a reference. Without a specific pressure set it is empty.}

\prop{2.14}{d}{A \textbf{class} is bounded by shared \textit{affordance structure}.\textsuperscript{12,\,55,\,56} The objects in the class offer the same operations to the same agents.}

\prop{2.141}{t}{The affordance structure is constant within the class. It carries zero information about why one form arose rather than another.}

\prop{2.142}{o}{Style predicts more form variation than material.\textsuperscript{53} Material is therefore not the driving force.}

\prop{2.2}{d}{\textbf{Material affordance} consists of the operations the material permits and prevents.\textsuperscript{12} Its signature is \textit{latent}; it binds the probability distribution without forcing the geometry.}

\prop{2.201}{o}{Steel existed for centuries before the tubular chair because the operation was missing.\textsuperscript{54} Form history is systematically slower than material history.}

\propcont{The tube lay in the bicycle factory for half a century. Someone had to see the chair in it.}

\prop{2.21}{t}{An \textit{aggregate variable} predicts form variation better than single pressures. The style period is such an aggregate variable.}

\prop{2.211}{t}{The style is not an explanation; it is an aggregate that absorbs all simultaneous pressures under one label.}

\prop{2.212}{t}{The style is a symptom of underlying forces, not a cause.}

\flowchapter{3 The selection pressures produce a landscape}{the selection pressures produce a landscape}

\mainprop{3}{}{The selection pressures span a landscape over the morphospace.}

\prop{3.1}{d}{The \textbf{adaptation landscape} is the vector sum of all active selection pressures.\textsuperscript{3,\,8,\,58} Each position holds a value; the degree of simultaneous satisfaction of all pressures.}

\prop{3.11}{t}{The landscape can be sounded only retrospectively. It yields local prediction, but no single form is guaranteed. In stasis the line of sight is long; at bifurcation it breaks.}

\prop{3.12}{t}{The adaptation function has many local maxima. A peak is a place where every deviation costs; valleys are the states forms shun.}

\prop{3.121}{t}{One universal peak is a theoretical exception.\textsuperscript{58} The empirical record shatters the postulate of the single perfect form.}

\prop{3.2}{d}{The stability of a peak corresponds to the steepness of its valley walls.\textsuperscript{8} Steep walls compel \textit{canalisation}. The form repels disturbances.}

\prop{3.21}{t}{Deep channels freeze dimensions; shallow ones accommodate fluctuation.}

\prop{3.3}{d}{A \textbf{style} is the \textit{swarming} of form around a shared centre of gravity.\textsuperscript{9,\,53} The style is a \textit{swarm}, not a category; its boundaries are necessarily fuzzy. Style periods are gradients, not islands.}

\propcont{No one has ever seen a style. One sees chairs that resemble each other.}

\prop{3.4}{t}{Form-making does not create the landscape's basic structure; that follows from the material and energetic conditions. But each realisation modifies the landscape locally through the extension of K.}

\prop{3.41}{i}{Isolated traditions navigate the same cleft without exchanging a word. Independent convergence confirms that the attractor is real.}

\flowchapter{4 The landscape is dynamic}{the landscape is dynamic}

\mainprop{4}{}{The landscape is dynamic.}

\prop{4.1}{a}{The selection pressures never rest. Peaks rise, sink, shift, split, merge. Optimal adaptation is local in time.}

\prop{4.11}{o}{A new technology shifts the boundary of the forbidden; a new market changes what the landscape rewards; a new material extends the space of operations.}

\prop{4.12}{t}{Gradual pressure shifts accumulate until the peak collapses; the discontinuity lies in the topology, not in the dynamics.}

\prop{4.2}{o}{Long periods of stasis break into abrupt shifts. Radiation into a new region, then convergence. Discontinuity is the signature of the dynamics, not a break from it.}

\prop{4.3}{t}{The landscape has memory.\textsuperscript{14} Each realised form alters the selection pressures for the next. All form-making is re-form-making; the agent always inherits a position. \textit{Path dependence} is a fundamental feature, not an anomaly.}

\prop{4.4}{o}{The morphospace expands cumulatively. The \textit{adjacent possible}\textsuperscript{13,\,59,\,60,\,79} is the C\msym{∖}K boundary; the set of forms conceivable but unknown as realisable. Material flows drive the landscape change.}

\prop{4.5}{o}{When one pressure dominates, the landscape collapses towards a single attractor.\textsuperscript{59} A single-ruling pressure does not explain form; it is the absence of explanation. Diversity mirrors the plurality of active pressures.}

\prop{4.51}{t}{A doctrine that reduces all selection pressures to one collapses the field to a single dimension. This is not rationality; it is systematic blindness.\textsuperscript{15,\,78}}

\prop{4.511}{t}{Forms that arise under a collapsed landscape fail under the real pressures to which the doctrine was blind.}

\prop{4.512}{t}{The repressed competencies cannot be reconstructed from the doctrine that repressed them.}

\flowchapter{5 There are agents}{there are agents}

\mainprop{5}{}{The world contains agents that respond to the landscape.}

\prop{5.1}{a}{Every class has at least one system that navigates the landscape via negative feedback.\textsuperscript{62}}

\prop{5.11}{d}{An \textbf{agent} is an operator defined by the triple: target state, distance metric, adjustment mechanism.\textsuperscript{4,\,7,\,21,\,82,\,83,\,84} Agency requires only organisation; neither consciousness, intention, nor nervous system.}

\prop{5.111}{t}{The definition is substrate-independent.\textsuperscript{88,\,89} Any system satisfying the conditions is an agent.}

\propcont{Formally. Agent A := (g, d, \msym{δ}). There exists a Lyapunov function V such that}

\begin{center}\small
\begin{tabular*}{\textwidth}{@{\extracolsep{\fill}}lp{35mm}@{}}
\toprule
V(g) = 0 & The target is the zero point \\[3pt]
V(x) > 0 for x \msym{≠} g & Everything else is distance \\[3pt]
V is non-increasing along x\textsubscript{n+1} = \msym{δ}(x\textsubscript{n}) & The system approaches the target \\
\bottomrule
\end{tabular*}
\end{center}

\prop{5.12}{o}{Agency is \textit{stratified}: the structure is stable, the resolution varies.}

\prop{5.121}{i}{The thermostat: one target, one dimension, one switch.}

\prop{5.122}{i}{The tree: several gradients, no central processor. Each branch navigates locally.}

\prop{5.123}{i}{The colony: millions of agents; the nest emerges in the field between them.}

\prop{5.124}{i}{The mammal: prediction, memory, planning.\textsuperscript{64,\,65} The representation stretches across time.}

\prop{5.125}{i}{The designer: the representation stretches beyond the agent's own needs.}

\begin{center}
\includegraphics[width=0.5\textwidth]{termostat-designar.png}
\end{center}

\prop{5.126}{t}{The designer navigates on behalf of others who are going to live with the form.}

\prop{5.13}{t}{The agent responds to the gradient in the landscape, not to the form.\textsuperscript{63,\,64} She does not invent the target; she discovers it.}

\prop{5.131}{t}{Navigation requires only the local gradient.\textsuperscript{63} Local competency produces global manifestation.}

\prop{5.2}{d}{The \textbf{cognitive light cone} is the subset of the morphospace the agent can represent and manipulate.\textsuperscript{21,\,85,\,90} Everything outside the light cone lacks causal existence for the agent. Light cones span from cellular millisecond response to decade-long consensus.}

\prop{5.21}{t}{The light cone dictates the resolution of the morphospace. Within it, the agent does not perceive forms but \textit{affordances}: the operations the configuration permits.\textsuperscript{12,\,55}}

\prop{5.211}{t}{Extending the light cone actualises latent affordances; narrowing it deactivates real ones.\textsuperscript{55,\,70}}

\prop{5.212}{d}{A form is in the light cone if there exists a policy such that the information loss between what the agent observes and what the agent predicts internally is below a threshold \(\varepsilon\).\textsuperscript{86}}

\prop{5.213}{t}{The information loss is a measure of the uncertainty introduced when the agent's internal model is used to represent the observations.}

\prop{5.3}{d}{Configurations closed under union, difference, and transformation constitute a \textit{shape algebra}.\textsuperscript{11,\,13,\,18}}

\prop{5.301}{d}{The \textbf{shape grammar} performs discrete transitions in this geometry. Rules fire on \textit{embedding}: the agent recognises a subshape within the whole shape. Input geometry is the sole premise.}

\propcont{A shape grammar is a triple SG = (S, R, \msym{ω}) where R = \{r\textsubscript{i} : (a\textsubscript{i}, b\textsubscript{i})\}. The rules transform shapes:}

\begin{center}\small
\begin{tabular*}{\textwidth}{@{\extracolsep{\fill}}p{22mm}lp{26mm}@{}}
\toprule
\textbf{Rule application} & s \msym{→} (s \msym{−} \msym{τ}(a\textsubscript{i})) + \msym{τ}(b\textsubscript{i}) & Find the subshape; replace it \\[4pt]
\textbf{Language} & L(SG) = \{s \msym{∈} S : \msym{ω} \msym{→}* s\} & All shapes the grammar can generate \\[4pt]
\textbf{Firing} & & The rule fires for every embedding of a\textsubscript{i} in s \\[4pt]
\textbf{Premise} & & The current geometry; not construction history \\
\bottomrule
\end{tabular*}
\end{center}

\prop{5.31}{t}{The grammar operation is a C\msym{→}C partition; it extends the conception space. Realisation is a C\msym{→}K operation; it decides the configuration.}

\prop{5.311}{t}{The grammar defines the space of possibility; the landscape fixes the trajectory.}

\prop{5.312}{t}{The agent selects a configuration that reduces the local gradient.}

\prop{5.4}{d}{The substrate functions as a \textit{zeroth-order agent}.\textsuperscript{26,\,81,\,92,\,93} It vetoes configurations and gives resonance to others through inexorable constraints.}

\prop{5.401}{t}{The substrate has two layers: structure parameters (the initial configuration) and function parameters (the competency rules operating on the configuration).\textsuperscript{92} Under noise, function parameters accelerate navigation in the morphospace.}

\prop{5.41}{i}{Biological tissue navigates via bioelectric polarisation;\textsuperscript{22,\,24,\,85} human practice via material friction; a slime mould via oscillating contractions; a language model via vectorial attention.}

\prop{5.5}{t}{Navigation competency accumulates as structured information.\textsuperscript{64} Learning is the irreversible expansion of accessible morphospace.}

\prop{5.51}{t}{Competency accumulation extends C(A) but does not guarantee care; the agent may know more axes without adjusting for them.}

\prop{5.6}{d}{\textbf{Care}\textsuperscript{23,\,68} is the set of axes in the morphospace the agent \textit{actively} represents and adjusts for.}

\prop{5.601}{t}{Care is not capacity. Capacity is how many axes the agent \textit{can} represent; care is how many it actually represents.}

\prop{5.61}{t}{Care determines how many axes the light cone contains. More axes reveal more compromises.}

\prop{5.62}{d}{Competency is the capacity to navigate towards stable positions in the landscape. Competency grows with care. Intelligence is this competency; care is the content of intelligence.\textsuperscript{23,\,86,\,91}}

\prop{5.621}{d}{Competency is quantifiable. Search efficiency \(K = \log_{10}(\tau_{\mathrm{blind}} / \tau_{\mathrm{agent}})\) is the number of orders of magnitude of work the agent's policy saves relative to a random walk.\textsuperscript{86} \(K = 0\) is the absence of navigation; \(K > 0\) is competency.}

\prop{5.63}{t}{An agent without care for a dimension is blind to the gradient along that axis. The blindness is not neutral; the invisible pressures still act.}

\prop{5.631}{t}{Care cannot be supplied afterwards; it had to be there when the form was navigated.}

\prop{5.65}{d}{\textbf{Life} is the process in which a system maintains and extends its own light cone over time.\textsuperscript{85,\,89}}

\prop{5.651}{t}{A system is living to the degree that it actively counteracts reduction of its light cone.}

\prop{5.652}{t}{Death is the irreversible collapse of the light cone to the empty.}

\prop{5.653}{t}{The boundary between living and non-living is not ontological.\textsuperscript{89} All persistent systems minimise variational free energy and are formally Bayesian actors. The difference is in degree, not in kind.}

\prop{5.7}{d}{\textbf{The core equation.}\textsuperscript{11,\,17,\,21} Form is the grammar's response to the gradient in the landscape, projected onto the agent's light cone:}

\begin{center}\small
\begin{tabular*}{\textwidth}{@{\extracolsep{\fill}}lp{58mm}@{}}
\toprule
Form = SG(\,\msym{∇}\textsubscript{C(A)}\, \msym{L}(\textit{c},\,\textit{t}\,)) & Grammar, landscape, light cone \\
\bottomrule
\end{tabular*}
\end{center}

\prop{5.71}{o}{The equation is a structural schema. It names the components that are necessary and together sufficient. It does not specify the functional form of \msym{L}, which rules in R fire, or the dimensionality of C(A).}

\prop{5.72}{t}{\textbf{Reduction.}\textsuperscript{11,\,17} The equation reduces to two limiting cases:}

\propcont{\textit{(i) Pure shape grammar}, when \msym{∇}\msym{L} \msym{≡} 0, |A| = 1, and \msym{Π} = (M, \msym{∅}, \msym{∅}). Then Form = SG(0) = L(SG): the full generated language without selection.}

\propcont{\textit{(ii) Pure concept--knowledge theory}, when SG = id\textsubscript{M} and \msym{L} = \msym{𝟙}\textsubscript{K}. The operations then reduce to partition changes between C and K (the operators in 1.443).}

\prop{5.721}{t}{Shape grammar and concept--knowledge theory are therefore not competing frameworks. They are degenerate limiting cases of the same equation; each describing the situation in which one of three components is trivial.}

\prop{5.722}{t}{Since historical design situations are never degenerate, each of the two frameworks alone is systematically insufficient. They do not lack details; they lack a layer.}

\prop{5.723}{a}{The equation closes the field. More layers cannot be added without changing the structure of the world of forms itself; fewer cannot be removed without disappearing into a limiting case.}

\flowchapter{6 Form arises between the agents}{form arises between the agents}

\mainprop{6}{}{Form arises in the field between agents.}

\prop{6.01}{d}{A \textbf{form-maker} is an agent that navigates the morphospace. The cell, the bee, the tree, the craftsperson, and the engineer are form-makers; it is the substrate of the agency that varies, not its structure.}

\prop{6.02}{d}{A \textbf{designer} is a form-maker whose target state includes the objectives of agents that themselves have no access to the navigation space. The differentia specifica is the entrusted trust; the designer navigates on behalf of others who will live with the form.}

\prop{6.021}{t}{All designers are form-makers. Not all form-makers are designers. The cell navigates on its own behalf within its own light cone; it does not represent its future users because it has none. The designer does precisely that.}

\prop{6.022}{t}{The designer has the capacity to refuse the navigation. Other form-makers have no such capacity; the cell builds what it builds, the tree grows where it grows. Responsibility follows from the asymmetry of choice, not from an asymmetry of insight.}

\prop{6.1}{t}{Several agents respond to different pressures simultaneously. Each form is a \textit{poly-agential compromise}.\textsuperscript{66,\,67} No single agent dictates the morphology; it emerges at the intersection.}

\prop{6.11}{t}{Each level solves problems competently within its own light cone.\textsuperscript{67} Lower components promote targets they lack the capacity to represent.}

\prop{6.111}{t}{Reductionism and holism are complementary.\textsuperscript{72} Both are true; neither is sufficient alone.}

\prop{6.12}{t}{Agency is \textit{multiscalar}.\textsuperscript{22,\,67,\,85,\,87} Each scale has its own competency and its own light cone. Form production is distributed C\msym{↔}K expansion carried out by grammar operations across competency scales.}

\prop{6.121}{t}{Evolutionary transitions in individuality are the biological mechanism for scale jumps.\textsuperscript{87} Deep interaction structure; not symmetric interplay; is the condition for a system of agents to become a single agent.}

\begin{center}
\includegraphics[width=0.9\linewidth]{formgjevarkompetanse/agenthierarki.png}
\end{center}

\prop{6.13}{t}{The agent chooses; the landscape filters.\textsuperscript{66} Control is always distributed. Exploration requires loose couplings; exploitation requires tight ones.}

\prop{6.14}{t}{The designer is never alone in the landscape.\textsuperscript{66,\,67,\,93} She navigates in a field of competing agencies from microbial tissue to market forces. Her competency lies in coordinating independent light cones into a stable form-optimum.}

\prop{6.2}{o}{Tight coupling establishes global feedback loops; the whole itself becomes an agent.\textsuperscript{66}}

\prop{6.3}{t}{A \textit{style} lacks causal power;\textsuperscript{9,\,15,\,53} it is not a force, but a \textit{pattern memory}. The style functions as a macroscale attractor that canalises the light cones of the agents who observe it.}

\propcont{It simplifies navigation by offering a recognisable centre of gravity; but to design ``in a style'' without understanding the pressures that generated the cluster is a category mistake.}

\prop{6.31}{o}{For the user, the style is not a symptom; it is an interface.\textsuperscript{80} The user perceives the centre of gravity, not the selection pressures.}

\prop{6.311}{t}{The designer who ignores this ignores the agent she is designing for.\textsuperscript{80}}

\prop{6.32}{t}{To navigate with style knowledge is to represent more axes.}

\prop{6.4}{t}{Each realised form extends the known and shifts the boundary of the adjacent possible (cf. 4.4).\textsuperscript{14,\,79} The history of form does not exhaust a closed repertoire; it extends it through movement.}

\prop{6.5}{t}{No form is final. Each balance becomes suboptimal in time.\textsuperscript{60,\,61} Only the generative structure endures; morphospace, pressure, landscape, agents. The forms are fleeting traces; the grammar is stable.}

\flowchapter{7 The form shows how far care reached}{the form shows how far care reached}
\mainprop{7}{}{The form shows how far care reached.}

\prop{7.1}{t}{The realised form is a legible archive of the compromise's dimensionality. A form designed under narrow care fails systematically along the axes that were not attended to. The failure is not accidental; it carries the signature of the blindness that produced it.}

\prop{7.11}{t}{The signature is legible because form is not random noise; it is the residue of a directed process. The direction is given by the light cone.}

\prop{7.12}{t}{Two forms designed under different care for the same function will differ along the axes the one designer saw and the other did not see. The difference is the archive.}

\prop{7.13}{t}{The archive is not an accusation. It is a fact. The same fact that makes geology legible makes design legible.\textsuperscript{72}}


\prop{7.15}{t}{A design theory without a concept of cognitive light cone cannot predict failure signature.\textsuperscript{21,\,70} It can only register it after the fact.}

\prop{7.151}{t}{The prediction requires that the theory names the representational boundary.\textsuperscript{70} Without the boundary there is no axis to single out as unrepresented before failure occurs.}

\prop{7.152}{t}{A shape grammar alone (5.72 i) generates the possible language; it cannot distinguish what in the language will fail along what axis. A concept--knowledge theory alone (5.72 ii) partitions K and C; it cannot derive failure signature. Prediction therefore requires all three layers.}

\prop{7.153}{t}{The failure signature is therefore not an empirical bonus; it is the touchstone that the theory actually explains form. A theory that does not predict failure signature does not explain form; it catalogues it.}

\prop{7.2}{a}{To extend the light cone is an ethical act.\textsuperscript{68,\,90} To narrow it is to surrender parts of the world of forms to forces no one represents.}

\prop{7.201}{t}{The consequences are borne by those who live with the form, not by the one who gave it.}

\prop{7.202}{a}{If the designer could have represented more axes but chose not to, the resulting failures are attributable to that choice.\textsuperscript{68,\,69}}

\prop{7.21}{t}{The act of design has real consequences for other agents navigating the same landscape.\textsuperscript{68}}

\prop{7.22}{a}{Consequences that an agent had the capacity to see but chose not to represent are attributable to that choice.}

\prop{7.23}{t}{The gap between capacity and care is the ethical space in design.\textsuperscript{68} Capacity sets the upper bound; care decides how much of the bound is used. The unused gap is what the form fails along.}

\prop{7.3}{t}{The world of forms is the case that acts. Navigation continues.}

\flowchapter{8 The generative model}{the generative model}

\mainprop{8}{}{Form is a probability distribution over the morphospace.}

\prop{8.1}{d}{The \textbf{generative model} for a form F realised by agents \{A\textsubscript{i}\} navigating in a landscape L with grammar SG is:\textsuperscript{86,\,89}}

\begin{center}\small
\begin{tabular*}{\textwidth}{@{\extracolsep{\fill}}lp{36mm}@{}}
\toprule
\(P(F \mid \{A_i\}, L, \mathrm{SG}) \propto \exp\!\bigl(-\textstyle\sum_i w_i \hat{L}_i(F)\bigr)\) & Boltzmann weighting \\[4pt]
\(F \in L(\mathrm{SG}) \cap \bigl(\bigcap_i C(A_i)\bigr)\) & The domain \\
\bottomrule
\end{tabular*}
\end{center}

\prop{8.11}{t}{\(\hat{L}_i\) is agent \(i\)'s estimate of the adaptation landscape. \(w_i \geq 0\) is the agent's relative influence with \(\sum_i w_i = 1\).}

\prop{8.12}{t}{The proportionality constant follows from the requirement that P sums to one over the domain \(L(\mathrm{SG}) \cap (\bigcap_i C(A_i))\).}

\prop{8.2}{t}{The model unites three components: grammar (the space of possibility \(L(\mathrm{SG})\)), landscape (the preference \(\hat{L}_i\)), and light cone (the representation \(C(A_i)\)).}

\prop{8.21}{t}{The three components are not independent. The light cone constrains which part of the grammar's language is accessible. The landscape constrains which part of the accessible is probable.}

\prop{8.3}{t}{The limiting cases of the model are:\textsuperscript{92}}

\propcont{(i) Flat landscape: \(\hat{L}_i(F) \equiv 0\) for all F and i. All forms in \(L(\mathrm{SG}) \cap (\bigcap_i C(A_i))\) are equally probable. Pure shape grammar without selection.}

\propcont{(ii) Trivial grammar: \(L(\mathrm{SG}) = \{\omega\}\). Only the start state is in the language. Pure gradient navigation in a fixed space.}

\prop{8.31}{t}{Formlære is not an alternative to shape grammar theory or to concept--knowledge theory. It is a unification that shows their place as degenerate limiting cases of proposition 8.1.}

\begin{center}
\includegraphics[width=0.92\linewidth]{formgjevarkompetanse/den-generelle-forma.png}
\end{center}

